\section*{Chapitre \ref{ch:qm:markov-chains-and-queues} --- Chaînes de Markov et files d'attente}

\begin{solution}{exo:qm:markov-chains-and-queues:1}
Durées moyennes de séjour $1/q_i$~: 0.5~s à un tick, $1/6$~s à deux, $2/9 = 0.22$~s à trois. Depuis trois ticks, la chaîne passe à
deux avec la probabilité $4/4.5 = 0.889$ et à un avec $0.111$.
\end{solution}

\begin{solution}{exo:qm:markov-chains-and-queues:2}
$\rho = 0.8$~: $L = 0.8/0.2 = 4$ messages, $W = 1/(1 - 0.8) = 5$ microsecondes, et $L = \lambda W$.
\end{solution}

\begin{solution}{exo:qm:markov-chains-and-queues:3}
$(1 - (1/0.9)^{20})/(1 - (1/0.9)^{40}) = 10.8\%$.
\end{solution}

\begin{solution}{exo:qm:markov-chains-and-queues:4}
Non. $\pi_1q_{12} = 0.676 \times 2 = 1.35$ mais $\pi_2q_{21} = 0.265 \times 5 = 1.32$~: l'\omterm{def:qm:markov-chains-and-queues:balance}{équilibre détaillé} n'est pas vérifié. Le critère de Kolmogorov
compare $q_{12}q_{23}q_{31} = 2 \times 1 \times 0.5 = 1$ à $q_{13}q_{32}q_{21} = 0$~: la probabilité circule $1 \to 2 \to 3 \to 1$.
\end{solution}

\begin{solution}{exo:qm:markov-chains-and-queues:5}
Moyenne $20/(1 - 0.9) = 200$~s, médiane 137~s, et 15.6~\% en moins d'une minute. La loi est asymétrique à droite~: la plupart des trajectoires
se vident plus tôt que la moyenne~: quelques-unes s'écartent d'abord loin vers le haut et mettent très longtemps.
\end{solution}

\begin{solution}{exo:qm:markov-chains-and-queues:6}
D'après la loi de Little, $W = L/\lambda = 400/50 = 8$~secondes.
\end{solution}

\begin{solution}{exo:qm:markov-chains-and-queues:7}
73.4~\% sur 20\,000 courses simulées à graine fixée, contre 73.2~\% par uniformisation~: l'écart reste dans l'erreur type
de la simulation, égale à 0.3 point.
\end{solution}

\begin{solution}{exo:qm:markov-chains-and-queues:8}
Le 95~\% est la probabilité que le prochain tick soit à la hausse, et non un taux d'exécution~: un ordre placé au fond du gros cours acheteur s'exécute
avant que le faible cours vendeur ne se vide que dans 73~\% des cas, et moins s'il est plus loin dans la file. Et les exécutions ne sont pas un
échantillon aléatoire~: un ordre au cours acheteur s'exécute quand les vendeurs continuent de le frapper, souvent juste avant que le prix ne baisse. La
décision requiert le mark-out après les exécutions, et non le taux d'exécution.
\end{solution}

\begin{solution}{pb:qm:markov-chains-and-queues:1}
\textbf{1.} 80 ordres devant~; 20 au cours vendeur.
\textbf{2.} $80/1 + 1/0.6 = 81.7$~secondes.
\textbf{3.} Moyenne $20/(1 - 0.9) = 200$~s~; médiane 137~s.
\textbf{4.} 15.6~\%.
\textbf{5.} Une seule~: la file dérive vers le bas à 0.1 ordre par seconde (99.6~\% des trajectoires en moins de 20 minutes).
\textbf{6.} 73.2~\%.
\textbf{7.} 73.4~\%.
\textbf{8.} 99.0~\% avec 40 ordres au cours vendeur~; avec 80, le cours vendeur ne se vide pratiquement jamais en premier.
\textbf{9.} 99.3~\%.
\textbf{10.} 95.4~\% à la hausse avec 80 contre 20~; 50~\% avec des files égales.
\textbf{11.} 10.8~\%.
\textbf{12.} 156.6~secondes.
\textbf{13.} Les départs croîtraient avec la file, tirant vers le bas les files longues et vers le haut les courtes~: une taille à \omterm{def:qm:stochastic-differential-equations:ou}{retour à la moyenne}
plutôt qu'une dérive constante, et moins de files extrêmes.
\textbf{14.} Les rafales vident certaines files bien plus vite, et d'autres durent plus longtemps~: des queues de distribution plus épaisses des deux côtés, et une probabilité
d'exécution qui dépend de l'état du flux d'ordres.
\textbf{15.} La dépendance entre les deux côtés~: un gros achat vide le cours vendeur et fait arriver de nouveaux cours acheteurs, et les fournisseurs de
liquidité réagissent au déséquilibre (les modèles à files réactives du livre~10).
\textbf{16.} La différence de probabilité d'exécution entre le devant et le fond de la file, $100\%$ contre $73\%$ ici, multipliée par la
valeur d'une exécution (la moitié de la fourchette, diminuée du mouvement adverse attendu).
\textbf{17.} Le fait que le cours acheteur ait été frappé 80 fois avant que le cours vendeur ne se vide~: les vendeurs étaient agressifs, et le prix a plus de chances
de baisser ensuite que ne le suggèrent les chiffres inconditionnels.
\textbf{18.} Seulement si la moitié de la fourchette dépasse le mouvement adverse attendu après une exécution (le mark-out, livre~2,
chapitre~15)~; le modèle donne la probabilité d'exécution, pas la réponse.
\textbf{19.} Résultat nommé~: \emph{le fond de la file}. Avec 80 ordres devant et 20 au cours vendeur, l'ordre s'exécute avant
que le cours vendeur ne se vide avec une probabilité de 73~\% (73.2~\% exacte, 73.4~\% simulée), au bout de 82 secondes en moyenne~; avec un
cours vendeur deux fois plus long, 99~\%.
\textbf{20.} La course entre le départ des ordres devant et la vidange de la file de l'autre côté.
\end{solution}

\begin{solution}{iq:qm:markov-chains-and-queues:1}
Traitée comme une \omterm{def:qm:markov-chains-and-queues:mm1}{file M/M/1} avec $\rho = 0.9$~: 9 requêtes dans le système en moyenne (8.1 en attente), chacune y passant $1/(10\,000 - 9\,000)$
s $= 1$~ms, dix fois le temps de service.
\iqlookfor{$\rho/(1-\rho)$, loi de Little et explosion au voisinage de la saturation.}
\end{solution}

\begin{solution}{iq:qm:markov-chains-and-queues:2}
La matrice des taux de saut $q_{ij}$, avec pour diagonale l'opposé des sommes par ligne~; $P(t) = e^{t\mathcal L}$. La \omterm{def:qm:stochastic-differential-equations:kolmogorov}{loi stationnaire}
est solution de $\pi\mathcal L = 0$ avec $\sum\pi_i = 1$~; pour un \omterm{def:qm:markov-chains-and-queues:bd}{processus de naissance et de mort}, l'\omterm{def:qm:markov-chains-and-queues:balance}{équilibre détaillé} la donne sous forme close.
\iqlookfor{taux, durées de séjour et $\pi\mathcal L = 0$.}
\end{solution}

\begin{solution}{iq:qm:markov-chains-and-queues:3}
À la hausse~: le faible cours vendeur se vide plus tôt que le gros cours acheteur, et la première file à se vider fait bouger le prix. Avec les
taux du chapitre, 95~\%.
\iqlookfor{épuisement d'une file vu comme une course.}
\end{solution}

\begin{solution}{iq:qm:markov-chains-and-queues:4}
$3/10$, par arrêt optionnel~; le nombre moyen de pas est $3 \times 7 = 21$, à partir de $W_n^2 - n$.
\iqlookfor{ruine du joueur et $i(n - i)$.}
\end{solution}

\begin{solution}{iq:qm:markov-chains-and-queues:5}
Le devant de la file s'exécute en premier et plus souvent, avant que le prix ne s'éloigne, et le fond s'exécute surtout lorsque la file est
consommée par le flux d'ordres informé. On valorise l'ordre comme la différence de probabilité d'exécution multipliée par l'edge par exécution, avec des
probabilités d'exécution issues d'un modèle de file ou des données d'exécution propres à la firme, par position dans la file.
\iqlookfor{probabilité d'exécution selon la position, et sélection adverse.}
\end{solution}

\begin{solution}{iq:qm:markov-chains-and-queues:6}
Uniformisation~: choisir $L$ au moins égal à chaque taux de sortie, appliquer la matrice des sauts $I + \mathcal L/L$ à un vecteur (produits
matrice-vecteur creux) et pondérer les itérés par des probabilités de Poisson. Chaque terme est positif ou nul, donc il n'y a
pas d'annulation, contrairement à un développement de Taylor de $e^{t\mathcal L}$.
\iqlookfor{uniformisation, parcimonie et positivité.}
\end{solution}
